\documentclass[border=3pt]{standalone}
\usepackage[T1]{fontenc}
\usepackage[utf8]{inputenc}
\usepackage{lmodern}
\usepackage{amsmath,amssymb}
\usepackage{xcolor}
\usepackage{tikz}
\usepackage{fontawesome5}
\usepackage{ragged2e}
\usetikzlibrary{shapes,shapes.geometric,shapes.misc,arrows.meta,positioning,
                calc,fit,backgrounds,decorations.pathreplacing,
                decorations.markings,patterns,chains}
\usetikzlibrary{decorations.pathreplacing,decorations.markings}
\definecolor{primary}{HTML}{1A5276}
\definecolor{secondary}{HTML}{2E86C1}
\definecolor{accent}{HTML}{E74C3C}
\definecolor{lightbg}{HTML}{EBF5FB}
\definecolor{defbg}{HTML}{FDF2E9}
\definecolor{greenhl}{HTML}{27AE60}
\definecolor{purplehl}{HTML}{8E44AD}
\definecolor{goldhl}{HTML}{B7950B}
\definecolor{tealhl}{HTML}{117A65}
\definecolor{codebg}{HTML}{F4F6F7}
\definecolor{neutral}{HTML}{707B7C}
\definecolor{dom1}{HTML}{1F618D}
\definecolor{dom2}{HTML}{117864}
\definecolor{dom3}{HTML}{6C3483}
\definecolor{dom4}{HTML}{922B21}
\definecolor{dom5}{HTML}{B7950B}
% --- carried over from ccna.sty so figures match the book exactly
\tikzset{
  dev/.style={rounded corners=3pt, draw=#1, fill=#1!8, text=#1!45!black,
              line width=0.9pt, minimum width=1.45cm, minimum height=1.05cm,
              align=center, font=\scriptsize\bfseries, inner sep=2pt},
  wirestyle/.style={line width=0.9pt, neutral!75},
  xconn/.style={line width=0.9pt, neutral!75, dash pattern=on 2.5pt off 2pt},
}

\newcommand{\Rtr}[3]{\node[dev=dom3] (#1) at #2 {\faIcon{route}\\#3};}
\newcommand{\Sw}[3]{\node[dev=dom2] (#1) at #2 {\faIcon{sitemap}\\#3};}
\newcommand{\Msw}[3]{\node[dev=dom3] (#1) at #2 {\faIcon{layer-group}\\#3};}
\newcommand{\Host}[3]{\node[dev=neutral] (#1) at #2 {\faIcon{desktop}\\#3};}
\newcommand{\Lap}[3]{\node[dev=neutral] (#1) at #2 {\faIcon{laptop}\\#3};}
\newcommand{\Srv}[3]{\node[dev=dom1] (#1) at #2 {\faIcon{server}\\#3};}
\newcommand{\Ap}[3]{\node[dev=dom5] (#1) at #2 {\faIcon{wifi}\\#3};}
\newcommand{\Wlc}[3]{\node[dev=dom5] (#1) at #2 {\faIcon{broadcast-tower}\\#3};}
\newcommand{\Fw}[3]{\node[dev=dom4] (#1) at #2 {\faIcon{shield-alt}\\#3};}
\newcommand{\Cld}[3]{\node[dev=neutral] (#1) at #2 {\faIcon{cloud}\\#3};}
\newcommand{\Phone}[3]{\node[dev=dom5] (#1) at #2 {\faIcon{phone-alt}\\#3};}
\newcommand{\Prn}[3]{\node[dev=neutral] (#1) at #2 {\faIcon{print}\\#3};}

% A link. The optional argument takes tikz options (e.g. xconn for a
% trunk drawn dashed); the fourth argument labels the link.
\newcommand{\wire}[4][wirestyle]{%
  \draw[#1] (#2) -- node[font=\tiny, text=neutral!85!black, fill=white,
                          inner sep=1.2pt, midway] {#4} (#3);}
\newcommand{\plainwire}[3][wirestyle]{\draw[#1] (#2) -- (#3);}


\newenvironment{topology}[1][]
  {\begin{center}\begin{tikzpicture}[font=\scriptsize,#1]}
  {\end{tikzpicture}\end{center}}


\newcommand{\cmd}[1]{\texttt{\small #1}}
\newcommand{\term}[1]{\textbf{\textcolor{primary}{#1}}}
\newcommand{\chref}[1]{Chapter~\ref{ch:#1}}
\newcommand{\ccnafitw}[1]{#1}
\providecommand{\thefigure}{}
% --- progressive builds
\newcount\buildlevel \buildlevel=99
% \stepvis{n}{..} appears from level n onward; \steponly{n}{..} at
% level n alone, which is how a caption is swapped rather than
% accumulated as the build advances. \stepdim{n}{..} is the
% complement of \stepvis: it shows only BEFORE level n, so a
% pair of them draws an element faint and then solid. That keeps
% the canvas full from step 1, which matters for a figure whose
% shape is itself the lesson -- a five-rung ladder revealed one
% rung at a time leaves a void where the method should be.
\newcommand{\stepvis}[2]{\ifnum#1>\buildlevel\relax\else#2\fi}
\newcommand{\steponly}[2]{\ifnum#1=\buildlevel #2\fi}
\newcommand{\stepdim}[2]{\ifnum#1>\buildlevel #2\fi}
% Slide text does not hyphenate. A projected caption broken as
% ``uni-cast'' costs the reader a beat that print can afford and a
% lecture cannot, and the ragged edge it avoids is invisible at
% four metres anyway.
\hyphenpenalty=10000 \exhyphenpenalty=10000
\begin{document}
\buildlevel=2
% Slide build of Figure 20.1 -- DORA, one message per step.
%
% The book prints all four messages at once, which is right on paper: a reader
% wants the whole exchange in one glance and will re-read it. In a lecture the
% four-message figure is a wall the room copies down instead of watching. Here
% each message lands on its own beat, and the caption under it says what has
% just changed -- so by step 4 the class has watched a handshake happen rather
% than read a diagram of one.
%
% The colour carries the exam point: three arrows are dom4 (broadcast), the
% ACK alone is dom1 (unicast). Nothing is dimmed and nothing moves, so Morph
% simply fades each new arrow in over a fixed background.
\begin{tikzpicture}[font=\scriptsize]
  % Fixed canvas at every level, or Morph slides the lifelines sideways.
  \path (-0.7,-6.6) rectangle (10.9,0.6);

  \node[font=\scriptsize\bfseries, text=neutral!45!black] (cl) at (0,0)  {CLIENT};
  \node[font=\scriptsize\bfseries, text=neutral!45!black] (sv) at (10,0) {SERVER};
  \draw[neutral!50, line width=0.8pt] (0,-0.35) -- (0,-5.3);
  \draw[neutral!50, line width=0.8pt] (10,-0.35) -- (10,-5.3);

  \tikzset{
    bc/.style={-{Stealth[length=2.4mm]}, line width=1.1pt, dom4!80},
    uc/.style={-{Stealth[length=2.4mm]}, line width=1.1pt, dom1!80},
    lb/.style={font=\scriptsize\bfseries, above=1pt, midway},
    sb/.style={font=\tiny, text=neutral!85!black, below=1pt, midway},
    ghost/.style={neutral!22, line width=1.1pt},
    % A pinned bounding box crops rather than grows, so a caption without an
    % explicit width runs off the canvas and is silently cut in half.
    stepcap/.style={font=\bfseries, anchor=north west, text width=11.3cm,
                    align=left}}

  % Faint placeholders for the messages still to come. Four lines say "four
  % messages" from the first beat -- which the name DORA has already told the
  % room anyway -- while the labels still arrive one at a time. Without them
  % the canvas is three-quarters empty at step 1 and the class reads the void
  % as the end of the diagram.
  \stepdim{2}{\draw[ghost] (0,-2.1) -- (10,-2.1);}
  \stepdim{3}{\draw[ghost] (0,-3.2) -- (10,-3.2);}
  \stepdim{4}{\draw[ghost] (0,-4.3) -- (10,-4.3);}

  \stepvis{1}{%
    \draw[bc] (0,-1.0)  -- node[lb, text=dom4!45!black] {1. DHCPDISCOVER}
                           node[sb] {broadcast, src 0.0.0.0, dst 255.255.255.255} (10,-1.0);}
  \stepvis{2}{%
    \draw[bc] (10,-2.1) -- node[lb, text=dom4!45!black] {2. DHCPOFFER}
                           node[sb] {broadcast --- the client still has no address} (0,-2.1);}
  \stepvis{3}{%
    \draw[bc] (0,-3.2)  -- node[lb, text=dom4!45!black] {3. DHCPREQUEST}
                           node[sb] {broadcast, so other servers withdraw their offers} (10,-3.2);}
  \stepvis{4}{%
    \draw[uc] (10,-4.3) -- node[lb, text=dom1!45!black] {4. DHCPACK}
                           node[sb] {the lease is now a binding on the server} (0,-4.3);}

  % The binding only exists after the ACK. Showing it at step 4 is the whole
  % reason the sequence is worth animating: nothing is leased until then.
  \stepvis{4}{%
    \node[draw=dom1!60, fill=dom1!8, rounded corners=3pt, line width=0.8pt,
          font=\tiny\ttfamily, align=left, inner sep=5pt, anchor=north east]
        at (10.8,-4.75)
        {\textcolor{dom1!45!black}{binding}\\ 192.168.10.51\\ 08:00:27:aa:bb:cc};}

  \steponly{1}{\node[stepcap, text=dom4!45!black] at (-0.6,-5.6)
      {A host with no address shouts into the broadcast domain. Who can hear it?};}
  \steponly{2}{\node[stepcap, text=dom4!45!black] at (-0.6,-5.6)
      {The offer is a broadcast too --- the client has nothing to unicast to yet};}
  \steponly{3}{\node[stepcap, text=dom4!45!black] at (-0.6,-5.6)
      {Broadcast again, on purpose: the servers not chosen hear it and release};}
  \steponly{4}{\node[stepcap, text=dom1!45!black] at (-0.6,-5.6)
      {Only now is there a lease. Three of the four were broadcasts.};}
\end{tikzpicture}

\end{document}
